\documentclass[11pt,a4paper]{article}
\usepackage[T1]{fontenc}
\usepackage[utf8]{inputenc}
\usepackage{mathtools,amsthm}
\usepackage{newtxtext,newtxmath}
\usepackage[margin=25mm,headheight=15pt]{geometry}
\usepackage{microtype}
\usepackage{booktabs,array,longtable}
\usepackage[dvipsnames]{xcolor}
\usepackage[most]{tcolorbox}
\usepackage{enumitem}
\usepackage{fancyhdr}
\usepackage{titlesec}
\usepackage{hyperref}
\usepackage{aliascnt}
\usepackage[nameinlink,noabbrev]{cleveref}
\definecolor{OliveInk}{HTML}{47543E}
\definecolor{PaleOlive}{HTML}{F1F3EC}
\definecolor{WarmGray}{HTML}{66645D}
\hypersetup{colorlinks=true,linkcolor=OliveInk,citecolor=OliveInk,urlcolor=OliveInk,
 pdftitle={Sharp ordinal bounds for Hoare powerspaces},
 pdfauthor={ChatGPT, research manuscript prepared for Vladimir Reshetnikov},
 pdfsubject={Noetherian spaces, ordinal stature, Hoare powerspaces, finite-tail spectra}}
\titleformat{\section}{\Large\bfseries\color{OliveInk}}{\thesection}{0.7em}{}
\titleformat{\subsection}{\large\bfseries\color{OliveInk}}{\thesubsection}{0.7em}{}
\pagestyle{fancy}
\fancyhf{}
\fancyhead[L]{\small\color{WarmGray}Sharp ordinal bounds for Hoare powerspaces}
\fancyhead[R]{\small\color{WarmGray}Research manuscript}
\fancyfoot[C]{\thepage}
\renewcommand{\headrulewidth}{0.3pt}
\setlength{\parindent}{1.1em}
\setlength{\parskip}{3pt}
\setlength{\emergencystretch}{2em}
\setlist{itemsep=3pt,topsep=5pt}
\newtheorem{theorem}{Theorem}[section]
\newaliascnt{lemma}{theorem}
\newtheorem{lemma}[lemma]{Lemma}
\aliascntresetthe{lemma}
\crefname{lemma}{Lemma}{Lemmas}
\newaliascnt{proposition}{theorem}
\newtheorem{proposition}[proposition]{Proposition}
\aliascntresetthe{proposition}
\crefname{proposition}{Proposition}{Propositions}
\newaliascnt{corollary}{theorem}
\newtheorem{corollary}[corollary]{Corollary}
\aliascntresetthe{corollary}
\crefname{corollary}{Corollary}{Corollarys}
\theoremstyle{definition}
\newaliascnt{definition}{theorem}
\newtheorem{definition}[definition]{Definition}
\aliascntresetthe{definition}
\crefname{definition}{Definition}{Definitions}
\newaliascnt{example}{theorem}
\newtheorem{example}[example]{Example}
\aliascntresetthe{example}
\crefname{example}{Example}{Examples}
\newaliascnt{remark}{theorem}
\newtheorem{remark}[remark]{Remark}
\aliascntresetthe{remark}
\crefname{remark}{Remark}{Remarks}
\newcommand{\Cl}{\mathcal C}
\newcommand{\HH}{\mathsf H}
\newcommand{\Hfin}{\mathsf H_{\mathrm{fin}}}
\newcommand{\JJ}{\mathcal J}
\newcommand{\rk}{\operatorname{rk}}
\newcommand{\st}[1]{\lVert #1\rVert}
\newcommand{\cl}{\operatorname{cl}}
\newcommand{\Spec}{\operatorname{Spec}_{\HH}}
\newcommand{\NN}{\mathbb N}
\newcommand{\boxF}[1]{\mathord{\Box}(#1)}
\newcommand{\ds}{\mathbin{\sqcup}}
\newcommand{\natsum}{\mathop{\bigoplus}}
\newcommand{\natprod}{\mathop{\bigotimes}}
\newcommand{\E}{\mathcal E}
\numberwithin{equation}{section}

\begin{document}
\hypersetup{pageanchor=false}
\begin{titlepage}
\centering
\vspace*{14mm}
{\small\sffamily\color{OliveInk}ORDER THEORY \enspace/\enspace ORDINAL RANKS \enspace/\enspace NOETHERIAN SPACES\par}
\vspace{13mm}
{\LARGE\bfseries Sharp ordinal bounds\\[3pt]for Hoare powerspaces\par}
\vspace{7mm}
{\Large Extremal statures and the complete\\[2pt]$\boldsymbol{\omega+n}$ spectrum\par}
\vspace{13mm}
{\large A research manuscript prepared for Vladimir Reshetnikov\par}
\vspace{3mm}
{\normalsize ChatGPT\par}
\vspace{3mm}
{\normalsize 19 September 2026\par}
\vspace{13mm}
\begin{minipage}{0.91\textwidth}
\begin{abstract}
The stature of a Noetherian space is the ordinal rank of the top element
of its lattice of closed subsets.  We study the stature of its extended
Hoare powerspace, equipped with the lower Vietoris topology and including
the empty closed set.  Motivated by Question~1 in Section~15 of
Goubault-Larrecq and Laboureix, we give a proof of the universal bound
$1+\st X\leq\st{\HH X}\leq 2^{\st X}$, where exponentiation is ordinal.
The binary upper bound is already known for the finitary powersets of
well-quasi-orders; the argument here extends it to all Noetherian spaces.
Explicit $T_1$ examples attain the upper bound at every ordinal stature.
Consequently the older bound $\omega^\alpha$ is attainable at infinite
input stature $\alpha$ exactly when $\omega\alpha=\alpha$.

We also determine the entire spectrum at each stature $\omega+n$,
$n<\omega$.  Its elements are precisely
$\omega\,j(U)+j(Q)-j(U)$, where $Q$ ranges over $n$-element posets,
$U$ ranges over their downsets, and $j(P)$ counts the downsets of $P$.
The proof isolates a canonical infinite closed core, removes a possible
generic point without changing the closed lattice, and computes a
finite-cell rank.  Every permitted value is realized by an explicit
countable space.  An accompanying program enumerates the finite spectra
through $n=6$ and checks a separate finite-model identity.
\end{abstract}
\end{minipage}
\vfill
\begin{minipage}{0.91\textwidth}
\small\color{WarmGray}
\textbf{Status and scope.} This is an original research attempt with full
English proofs, not an independently refereed or proof-assistant-verified
paper.  The general spectrum beyond the finite tails of $\omega$ is not
classified here.  Priority is not asserted: the distinction between the
known wqo theorem and the proposed general extension is explicit.
\end{minipage}
\end{titlepage}
\hypersetup{pageanchor=true}
\pagenumbering{roman}

\tableofcontents
\clearpage
\pagenumbering{arabic}

\section{The chosen question and the results}
\label{sec:intro}

The problem selected for this investigation is a concrete ordinal-spectrum
question rather than a foundational consistency question.  It has two
advantages for an attempted solution: the invariant is defined by a
well-founded recursion, and the relevant hyperspace has a finite normal
form for its closed sets.

Goubault-Larrecq and Laboureix establish
\begin{equation}\label{eq:old}
  1+\st X\ \leq\ \st{\mathcal H_{0V}X}\ \leq\ \omega^{\st X}.
\end{equation}
Their Section~15, Question~1 asks which intermediate ordinals occur, and
whether the upper endpoint can occur for \emph{every} infinite value of
$\st X$ \cite[Section~15, Question~1]{GL23}.  We write $\HH X$ for their
$\mathcal H_{0V}X$ throughout.  In particular, our powerspace includes the
empty closed subset of $X$; omitting that point would change several finite
and successor formulas.

For an ordinal $\alpha$, define
\[
 \Spec(\alpha)=\{\st{\HH X}:X\text{ is Noetherian and }\st X=\alpha\}.
\]
All constructions take place in ordinary set theory with choice.  Spaces
may have arbitrary cardinality and need not satisfy a separation axiom
unless one is stated.

\begin{tcolorbox}[colback=PaleOlive,colframe=OliveInk,boxrule=0.5pt,
 arc=1mm,title={Main results proved in this manuscript},fonttitle=\bfseries]
For every ordinal $\alpha$,
\[
 \min\Spec(\alpha)=1+\alpha,
 \qquad \max\Spec(\alpha)=2^\alpha.
\]
The maximum has a $T_1$ witness.  For infinite $\alpha$,
\[
 \omega^\alpha\in\Spec(\alpha)
 \quad\Longleftrightarrow\quad \omega\cdot\alpha=\alpha.
\]
For $n<\omega$, with $j(P)$ the number of downsets of a finite poset $P$,
\[
 \Spec(\omega+n)=
 \bigl\{\omega\cdot j(U)+(j(Q)-j(U)):
          |Q|=n,\ U\in\JJ(Q)\bigr\}.
\]
Here $U$ carries the induced order, and the empty downset is counted.
\end{tcolorbox}

The simplest obstruction to universal upper attainment is already decisive: $\Spec(\omega)=\{\omega\}$, so $\omega^\omega$ is not attainable at
input stature $\omega$.  The stronger result is not just this counterexample,
but the optimal upper bound for every input and a criterion for exactly
when the old bound coincides with it.

\subsection{What is already known, and what is being extended}

Abriola, Halfon, Lopez, Schmitz, Schnoebelen, and Vialard prove
\[
  1+o(A)\leq o(\mathcal P_f(A))\leq 2^{o(A)}
\]
for well-quasi-orders, and provide sharpness witnesses for all ordinals
\cite[Theorems~3.3 and~4.3]{AHLLSSV24}.  Thus neither binary ordinal
exponentiation nor its sharpness in the wqo setting is a new discovery
here.  Our finite/limit/successor proof strategy is closely parallel to
that existing theorem.  The added step is a topological rank decomposition
that works without an Alexandroff topology or a wqo specialization order.

This distinction matters.  An infinite $T_1$ Noetherian space has equality
as its specialization order, an infinite antichain, so it does not come
from a wqo in this way.  The cofinite topology is the first example.  The
$T_1$ witnesses below therefore give a realization in a genuinely different
class of spaces, although unrestricted sharpness witnesses already follow
from the known wqo results.

The proofs of the upper bound and the $\omega+n$ spectrum are supplied
from rank and topology definitions.  A separate known result, the product
formula for statures, is explicitly imported only for the transfinite
$T_1$ sharpness construction.  The source question and these relevant
primary papers were checked in their public versions; this is not a
certification that no later or unpublished work contains the same
extension.

\section{Ranks, finite pieces, and closed--open decompositions}
\label{sec:rank}

\subsection{Definitions and elementary monotonicity}

\begin{definition}
A space $X$ is \emph{Noetherian} when its closed subsets satisfy the
descending chain condition.  Equivalently, the inclusion order on
$\Cl(X)$, its set of closed subsets, is well-founded.  For $F\in\Cl(X)$,
put
\[
 r_X(F)=\sup\{r_X(E)+1:E\in\Cl(X),\ E\subsetneq F\},
 \qquad \st X=r_X(X).
\]
The supremum of the empty family is $0$.
\end{definition}

Thus $\st\varnothing=0$.  If $F$ is closed in $X$, its closed subsets are
exactly the members of the interval $[\varnothing,F]$ in $\Cl(X)$.
Consequently
\begin{equation}\label{eq:closedrank}
 \st F=r_X(F).
\end{equation}
There is no extra $1$ in this identity: stature is the rank of the top
point, not the rank of the whole closed-set poset including that top.

For well-founded posets, a strictly increasing map never decreases point
ranks.  Indeed, this follows by induction on the source point, since its
strict predecessors map to strict predecessors of its image.  We use
this simple observation repeatedly.

\begin{lemma}\label{lem:monotone}
If $Y$ is a subspace of a Noetherian space $X$, then $Y$ is Noetherian and
$\st Y\leq\st X$.  If $f:X\to Z$ is continuous and onto, then $Z$ is
Noetherian and $\st Z\leq\st X$.
\end{lemma}
\begin{proof}
For the subspace assertion, send a closed $A\subseteq Y$ to
$\cl_X(A)$.  The equality $Y\cap\cl_X(A)=A$ shows that this is an order
embedding of closed-set lattices.  Its image is bounded by $X$, so the
rank assertion follows.  It also prohibits an infinite descending chain
of closed subsets of $Y$.

For the image assertion, inverse image under $f$ is an order embedding
$\Cl(Z)\hookrightarrow\Cl(X)$, because $f$ is onto.  It sends $Z$ to $X$.
Apply the same reasoning.
\end{proof}

The Kolmogorov, or $T_0$, quotient identifies points with identical
closures.  Every closed subset is saturated for this equivalence, so the
quotient has an isomorphic closed-set lattice.  Passing to this quotient
therefore preserves stature.  Later we will see directly that it also
preserves the homeomorphism type of the Hoare powerspace.

\subsection{Natural sums and finite unions}

The \emph{natural sum} $\alpha\oplus\beta$ adds the coefficients in a
common Cantor normal form.  The \emph{natural product}
$\alpha\otimes\beta$ multiplies the corresponding polynomials in $\omega$,
using natural sums for exponents.  These operations are commutative and
are different from ordinary ordinal addition and multiplication.

\begin{lemma}\label{lem:productrank}
For well-founded posets $P,Q$,
\[
 \rk_{P\times Q}(p,q)=\rk_P(p)\oplus\rk_Q(q).
\]
Consequently, for a finite topological disjoint sum,
\[
 \st{X_1\ds\cdots\ds X_k}=
 \st{X_1}\oplus\cdots\oplus\st{X_k}.
\]
\end{lemma}
\begin{proof}
Natural sum satisfies the recursion
\[
 \alpha\oplus\beta=
 \sup\bigl(\{(\alpha'\oplus\beta)+1:\alpha'<\alpha\}
 \cup\{(\alpha\oplus\beta')+1:\beta'<\beta\}\bigr).
\]
The ranks realized strictly below a point of rank $\alpha$ include every
ordinal below $\alpha$: a proof appears in \cref{lem:gaps} below and is
independent of this lemma.  Applying the displayed recursion and induction
on $(p,q)$ gives the product formula.  Predecessors that decrease both
coordinates do not change the supremum, since they lie below a predecessor
that decreases just one coordinate.

The lattice of closed subsets of a finite disjoint sum is the product
of the closed-set lattices of its summands.  Evaluate the formula at its
top element and iterate.
\end{proof}

\begin{corollary}\label{cor:union}
For finitely many Noetherian subspaces $Y_i$ whose union is $Y$,
\[
 \st Y\leq\natsum_i\st{Y_i}.
\]
The $Y_i$ need not be closed in $Y$.
\end{corollary}
\begin{proof}
The map from the finite topological disjoint sum of the $Y_i$ onto their
union is continuous.  Use \cref{lem:monotone,lem:productrank}.
\end{proof}

\begin{lemma}[No gaps in point ranks]\label{lem:gaps}
If $p$ has rank $\alpha$ in a well-founded poset and $\beta<\alpha$, there
is $q<p$ with rank exactly $\beta$.
\end{lemma}
\begin{proof}
Choose a minimal member $q$ of
$\{x:x\leq p,\ \rk(x)\geq\beta\}$.  Every strict predecessor of $q$ has
rank less than $\beta$, so the rank recursion gives $\rk(q)\leq\beta$.
The reverse inequality holds by the choice of the set.  Since
$\beta<\rk(p)$, $q\ne p$.
\end{proof}

\begin{lemma}\label{lem:finite}
A finite $T_0$ space on $m$ points has stature $m$.  Conversely, if a
$T_0$ space has finite stature $m$, it has exactly $m$ points.
\end{lemma}
\begin{proof}
The closed sets of a finite $T_0$ space are the downsets of its
specialization poset.  Strict inclusion increases cardinality, giving an
upper bound $m$ on the rank of the top.  Initial segments of a linear
extension give a chain of $m+1$ closed sets from $\varnothing$ to $X$,
so the rank is $m$.

An infinite $T_0$ space contains finite $T_0$ subspaces of every finite
cardinality.  By \cref{lem:monotone}, its stature is at least every finite
ordinal.  This excludes finite stature for an infinite $T_0$ space.
\end{proof}

\subsection{The decomposition that drives the successor step}

\begin{lemma}[Closed--open rank inequality]\label{lem:semiadd}
For a closed subset $F$ of a Noetherian space $X$, put $V=X\setminus F$.
Then
\begin{equation}\label{eq:semiadd}
 \st F+\st V\ \leq\ \st X\ \leq\ \st F\oplus\st V.
\end{equation}
The addition on the left is ordinary ordinal addition.
\end{lemma}
\begin{proof}
Closed subsets $B$ of $V$ correspond, by $B\mapsto F\cup B$, to the
interval $[F,X]$ in $\Cl(X)$.  To verify this, write $B=V\cap C$ with
$C$ closed in $X$; then $F\cup B=F\cup C$ is closed.

An induction on $B$ gives
\[
 r_X(F\cup B)\geq\st F+r_V(B).
\]
At $B=\varnothing$ this is equality.  At a nonempty $B$, use predecessors
$F\cup B'$ with $B'\subsetneq B$ and the continuity of ordinary addition
in its right argument.  Evaluating at $B=V$ proves the lower inequality.

For the upper inequality, the identity from the finer disjoint-sum
topology $F\ds V$ onto $X$ is continuous.  Apply the finite-union bound.
\end{proof}

\begin{corollary}[Finite tail decomposition]\label{cor:tail}
Let $X$ be $T_0$ with $\st X=\lambda+n$, where $\lambda$ is a nonzero
limit ordinal and $0<n<\omega$.  There exists a closed $F\subseteq X$
with $\st F=\lambda$, and its complement has exactly $n$ points.
\end{corollary}
\begin{proof}
Use \cref{lem:gaps} in the closed-set lattice to obtain $F$.  The lower
inequality in \eqref{eq:semiadd} gives
$\lambda+\st{X\setminus F}\leq\lambda+n$.  Ordinary addition is strictly
increasing in the right argument, so $\st{X\setminus F}\leq n$.
By \cref{lem:finite}, the complement has some finite cardinality $m\leq n$.
The upper inequality gives
\[
 \lambda+n\leq\lambda\oplus m=\lambda+m.
\]
Hence $m=n$.
\end{proof}

This is stronger than merely finding a small residual space: after
passing to $T_0$, a finite ordinal tail is an actual finite set of points.
Nothing here assumes that the specialization order is a wqo.

\section{The closed-set structure of a Hoare powerspace}
\label{sec:hoare}

\begin{definition}
For a space $X$, the points of $\HH X$ are its closed subsets, including
$\varnothing$.  Its \emph{lower Vietoris topology} has subbasis
\[
 \Diamond O=\{C\in\Cl(X):C\cap O\ne\varnothing\},\qquad O\text{ open in }X.
\]
For a closed $F\subseteq X$, write
\[
 \boxF F=\{C\in\Cl(X):C\subseteq F\}.
\]
\end{definition}

The set $\boxF F$ is closed, being the complement of
$\Diamond(X\setminus F)$.  Its points are themselves closed subsets of
$X$; this two-level notation should not be confused with $F$ as a
subspace of $X$.

\begin{proposition}\label{prop:closedH}
If $X$ is Noetherian, then $\HH X$ is Noetherian, and its closed subsets
are exactly
\begin{equation}\label{eq:closedH}
 D=\bigcup_{i=1}^k\boxF{F_i},\qquad F_i\in\Cl(X),\quad k<\omega.
\end{equation}
The case $k=0$ denotes the empty subset of the powerspace.
\end{proposition}
\begin{proof}
Finite unions in \eqref{eq:closedH} form a family $\mathcal D$ closed
under finite intersections, because
$\boxF F\cap\boxF G=\boxF{F\cap G}$.  They include the whole space.
We first show that $\mathcal D$ is well-founded under strict inclusion.

Each nonempty member has a unique finite antichain of maximal generators:
remove every $F_i$ contained in another generator.  For this irredundant
set of generators $M(D)$ define
\[
 w(D)=\natsum_{F\in M(D)}\omega^{r_X(F)},\qquad w(\varnothing)=0.
\]
Suppose $D\subsetneq D'$.  Each generator of $D$ lies below a generator
of $D'$.  A generator of $D$ not also a generator of $D'$ must lie
\emph{strictly} below some generator of $D'$ that is absent from $M(D)$;
otherwise the maximality of the generators of $D$ would fail.
There is at least one such removed generator.  Let $\gamma$ be the largest
rank of a generator removed from $M(D')$.  Every newly added generator
has rank below the rank of some removed generator, and hence below
$\gamma$.  At exponents above $\gamma$ the two Cantor normal forms agree;
at exponent $\gamma$ the coefficient decreases.  Thus $w(D)<w(D')$.

An arbitrary intersection of members of $\mathcal D$ is a finite
intersection.  Indeed, among all their finite intersections choose one
minimal under inclusion.  Intersecting it with any additional member
cannot make it smaller, by minimality.  It is therefore their full
intersection.  This proves that $\mathcal D$ is closed under arbitrary
intersections.  Since the closed subbasis consists of the sets
$\boxF F$, it follows that $\mathcal D$ is precisely the closed-set
family of $\HH X$.  Its well-foundedness proves Noetherianity.
\end{proof}

Several consequences are useful.  First, $\HH X$ is $T_0$, with
specialization order inclusion: the closure of a point $F$ is exactly
$\boxF F$.  Second, if $F$ is closed in $X$, then
\begin{equation}\label{eq:principalH}
 \boxF F\ \cong\ \HH F
\end{equation}
as spaces.  To check the topology, restrict each $\Diamond O$ to the
closed subsets of $F$ and replace $O$ by $O\cap F$.
Third, any isomorphism $\Cl(X)\cong\Cl(Y)$ induces a homeomorphism
$\HH X\cong\HH Y$, since \eqref{eq:closedH} depends only on that order.
This justifies the earlier reduction to the $T_0$ quotient.

\begin{lemma}\label{lem:functor}
If $f:X\to Y$ is continuous and onto, the map
\[
 f_*:\HH X\longrightarrow\HH Y,\qquad C\longmapsto\cl_Y(f[C])
\]
is continuous and onto.  For a finite disjoint sum,
\begin{equation}\label{eq:Hsum}
 \HH(X_1\ds\cdots\ds X_k)\cong\HH X_1\times\cdots\times\HH X_k.
\end{equation}
\end{lemma}
\begin{proof}
An open set meets $\cl_Y(f[C])$ if and only if it meets $f[C]$, so
$f_*^{-1}(\Diamond O)=\Diamond(f^{-1}(O))$.  For a closed $B\subseteq Y$,
choose $C=f^{-1}(B)$ to obtain $f_*(C)=B$.

For the disjoint-sum assertion, a closed set is a tuple of closed sets
in the summands.  Under this bijection a hit-open set becomes a union
of coordinate hit-open sets.  Conversely, a coordinate hit-open set is
obtained by using an open subset contained in that summand.  These
observations give equality of the topologies.
\end{proof}

\begin{lemma}[Finite discrete extension]\label{lem:discrete}
Let $D_m$ be a discrete space of $m$ points.  For every Noetherian $F$,
\[
 \st{\HH(F\ds D_m)}\leq\st{\HH F}\otimes 2^m.
\]
\end{lemma}
\begin{proof}
There are $2^m$ points in $\HH D_m$.  By \eqref{eq:Hsum}, the powerspace
in question is $\HH F\times\HH D_m$.  Cover it by its $2^m$ slices,
one for each point of the second factor.  Each slice is homeomorphic
to $\HH F$.  The finite-union bound gives the natural sum of $2^m$
copies of $\st{\HH F}$, which is $\st{\HH F}\otimes2^m$.
The slices need not be closed; \cref{cor:union} deliberately allows this.
\end{proof}

\subsection{Discarding non-finitary points without changing stature}

Let $\Hfin X$ be the subspace of $\HH X$ consisting of closures of finite
subsets of $X$.  A closure in this definition need not be finite as a set.
The following lattice observation is particularly important later; it is
the mechanism behind the equality also recorded in
\cite[Proposition~11.4]{GL23}.

\begin{proposition}\label{prop:finitary}
For Noetherian $X$, taking traces on $\Hfin X$ gives an isomorphism
\[
 \Cl(\HH X)\ \cong\ \Cl(\Hfin X).
\]
In particular, $\st{\HH X}=\st{\Hfin X}$.
\end{proposition}
\begin{proof}
Surjectivity is the definition of the subspace topology.  To prove order
reflection, suppose $D,E$ are closed in $\HH X$ and $D\nsubseteq E$.
Choose a point $C\in D\setminus E$ and write
$E=\bigcup_{i=1}^k\boxF{F_i}$.  For each $i$ choose
$x_i\in C\setminus F_i$.  The closed set
$A=\cl_X\{x_1,\ldots,x_k\}$ is contained in $C$, so it lies in $D$, but
it is contained in none of the $F_i$, so it lies outside $E$.
It is a point of $\Hfin X$.
If $k=0$, use $A=\varnothing$; every nonempty closed subset of $\HH X$
contains this point.  Thus the traces distinguish every failure of
inclusion.  An order isomorphism of closed lattices preserves top ranks.
\end{proof}

\section{Binary ordinal exponentiation}
\label{sec:ordinal}

All occurrences of $2^\alpha$ in this article mean ordinary
\emph{ordinal} exponentiation, not cardinal exponentiation and not a
natural exponentiation.  In particular $2^\omega=\omega$.

\begin{lemma}\label{lem:binary}
Write an ordinal uniquely as $\alpha=\omega\cdot q+n$, $n<\omega$.
Then
\begin{equation}\label{eq:binary}
 2^\alpha=\omega^q\cdot2^n.
\end{equation}
If $\lambda$ is a nonzero limit ordinal, $2^\lambda$ is a limit ordinal
and is indecomposable under finite natural sums of smaller ordinals.
Moreover, for all ordinals $\alpha,\beta$,
\begin{equation}\label{eq:hom}
 2^{\alpha\oplus\beta}=2^\alpha\otimes2^\beta.
\end{equation}
\end{lemma}
\begin{proof}
Ordinal exponent laws give
\[
 2^{\omega q+n}=(2^\omega)^q\cdot2^n=\omega^q\cdot2^n.
\]
For a nonzero limit, $n=0$ and $q>0$.  Every ordinal less than
$\omega^q$ has only Cantor normal form exponents less than $q$; a finite
natural sum still has this property.  This proves the indecomposability
assertion.

For \eqref{eq:hom}, left multiplication by $\omega$ changes each Cantor
normal form exponent $\xi$ to $1+\xi$.  This transformation is strictly
increasing, and hence preserves the merging of coefficients used in
natural sum.  Thus, writing $\alpha=\omega q+n$ and
$\beta=\omega r+m$,
\[
 \alpha\oplus\beta=\omega(q\oplus r)+(n+m).
\]
The finite remainders cannot combine with the terms having positive
exponent.  Formula \eqref{eq:binary} now yields
\[
 2^{\alpha\oplus\beta}
 =\omega^{q\oplus r}\,2^{n+m}
 =(\omega^q2^n)\otimes(\omega^r2^m).
\]
\end{proof}

The function $\alpha\mapsto2^\alpha$ is strictly increasing and continuous
at limit ordinals.  A limit input turns its value into a single
$\omega$-monomial.  That fact, rather than a coarse comparison with
$\omega^\alpha$, is the key to the next induction.

\section{The sharp universal upper bound}
\label{sec:upper}

\begin{theorem}\label{thm:upper}
For every Noetherian space $X$,
\begin{equation}\label{eq:mainbound}
 1+\st X\ \leq\ \st{\HH X}\ \leq\ 2^{\st X}.
\end{equation}
\end{theorem}
\begin{proof}
We may replace $X$ by its $T_0$ quotient.  Put $\alpha=\st X$ and prove
the upper bound by transfinite induction on $\alpha$.

\textbf{Finite input.}
If $\alpha=n<\omega$, then $X$ has exactly $n$ points.  The $T_0$ space
$\HH X$ is finite and has $|\Cl(X)|\leq2^n$ points.  Its stature equals
that cardinality, proving the bound.  This includes $X=\varnothing$,
whose powerspace is a singleton of stature $1=2^0$.

\textbf{Nonzero limit input.}
Suppose $\alpha$ is a limit ordinal.  Every proper closed
$D\subsetneq\HH X$ has the form
$\bigcup_{i=1}^k\boxF{F_i}$ with each $F_i\subsetneq X$.
Let $\beta_i=\st{F_i}<\alpha$.  By the finite-union bound, the principal
powerspace homeomorphism, and the induction hypothesis,
\[
 \st D\leq\natsum_{i=1}^k\st{\HH F_i}
       \leq\natsum_{i=1}^k2^{\beta_i}<2^\alpha.
\]
The final strict inequality is \cref{lem:binary}.  The empty $D$ satisfies
the same inequality.  Since $2^\alpha$ is a limit ordinal,
$\st D+1<2^\alpha$.  Taking the supremum over all proper closed $D$ gives
$\st{\HH X}\leq2^\alpha$.

\textbf{Infinite successor input.}
Write $\alpha=\lambda+n$, where $\lambda$ is a nonzero limit and
$0<n<\omega$.  By \cref{cor:tail}, choose a closed $F$ of stature
$\lambda$ with an $n$-point complement.  Replace the topology on this
complement by the discrete topology and separate it from $F$.  The
resulting identity $F\ds D_n\to X$ is continuous and onto, so
\cref{lem:functor,lem:discrete} give
\[
 \st{\HH X}\leq\st{\HH(F\ds D_n)}
 \leq\st{\HH F}\otimes2^n
 \leq2^\lambda\otimes2^n
 =2^{\lambda+n}.
\]
Here the penultimate step uses the induction hypothesis at $\lambda$;
the final identity follows from \cref{lem:binary}.

For the lower bound, $\boxF\varnothing=\{\varnothing\}$ has stature
$1$.  Induction on $F\in\Cl(X)$ gives
\[
 r_{\HH X}(\boxF F)\geq1+r_X(F).
\]
Indeed, each $\boxF E$ with $E\subsetneq F$ is a proper closed subset
of $\boxF F$, and ordinary left addition by $1$ is continuous in its
right argument.  The base $F=\varnothing$ is the singleton just noted.
Evaluate at $F=X$, for which $\boxF X=\HH X$.
\end{proof}

The upper-bound argument does not use the general product formula for
Noetherian spaces.  Its only product estimate is the finite-slice cover
in \cref{lem:discrete}.  In particular, the key extension beyond wqos is
not hidden inside an unproved product assertion.

\section{Sharpness at every ordinal, with \texorpdfstring{$T_1$}{T1} upper witnesses}
\label{sec:sharp}

\subsection{The lower endpoint}

\begin{proposition}\label{prop:lowerSharp}
For every ordinal $\alpha$, there is a Noetherian space $L_\alpha$ with
\[
 \st{L_\alpha}=\alpha,
 \qquad \st{\HH L_\alpha}=1+\alpha.
\]
\end{proposition}
\begin{proof}
Give the well-order $\alpha$ its Alexandroff topology of upward-closed
sets.  Its closed sets are its initial segments, of lengths
$\beta\leq\alpha$, so its stature is $\alpha$.

The finitary closed subsets are the empty set and the principal initial
segments $\mathord\downarrow x$, $x<\alpha$.  Under inclusion they form
a chain of type $1+\alpha$.  On this set of points the lower Vietoris
topology is exactly the Alexandroff topology: the upper cone of
$\mathord\downarrow x$ is the trace of $\Diamond(\mathord\uparrow x)$,
and the upper cone of the empty set is the whole space.  An ordinal
chain with this topology has stature its order type.  Therefore
$\st{\Hfin L_\alpha}=1+\alpha$; use \cref{prop:finitary}.
\end{proof}

\subsection{A finite-support bouquet}

We now use one substantial external theorem:
\begin{equation}\label{eq:productexternal}
 \st{Y\times Z}=\st Y\otimes\st Z
 \qquad (Y,Z\text{ Noetherian}).
\end{equation}
This is the product theorem of Goubault-Larrecq and Laboureix
\cite[Theorem~10.9]{GL23}.  The following all-ordinal $T_1$ sharpness
construction depends on it; the preceding upper bound and the later
spectrum classification do not.

\begin{definition}\label{def:bouquet}
For an infinite family $(X_i)_{i\in I}$ of nonempty Noetherian spaces,
let $B(X_i:i\in I)$ have underlying set $\coprod_{i\in I}X_i$.
Its closed subsets are the entire set and all sets
\[
 \coprod_{i\in S}F_i,
 \qquad S\subseteq I\text{ finite},\quad F_i\text{ closed in }X_i.
\]
We call this a \emph{finite-support bouquet}.  No additional generic
point is adjoined.
\end{definition}

\begin{lemma}\label{lem:bouquet}
The stated family is the closed-set family of a Noetherian topology.
Every finite union of whole components is closed and has the ordinary
finite-disjoint-sum topology.  If every $X_i$ is $T_1$, so is the bouquet.
\end{lemma}
\begin{proof}
Finite unions preserve finite support.  For an arbitrary intersection,
there is nothing to check if every member is the whole space.  Otherwise
fix one finite-support member; the intersection is confined to those
finitely many components and is closed in each of them.  It is therefore
again a permitted closed set.

Once a descending chain leaves the whole space, it lies in a finite
union of components.  This finite sum is Noetherian, so the chain
stabilizes.  The subspace-topology assertion follows directly by taking
traces of the defining closed sets.  Finally, every singleton is closed
when each summand is $T_1$.
\end{proof}

The finite-support idea is related to the successor construction in
\cite[Example~11.9]{GL23}.  Omitting an added generic point keeps the
present construction $T_1$, and the limit-stage construction below works
with an arbitrary ordinal-indexed family.

\begin{theorem}\label{thm:sharpT1}
For every ordinal $\alpha$ there is a $T_1$ Noetherian space $T_\alpha$
such that
\[
 \st{T_\alpha}=\alpha,
 \qquad \st{\HH T_\alpha}=2^\alpha.
\]
When $\alpha$ is countable, the witness can be countable.
\end{theorem}
\begin{proof}
First construct, by recursion on $\rho$, a space $B_\rho$ satisfying
\begin{equation}\label{eq:Btarget}
 \st{B_\rho}=\omega^\rho,
 \qquad\st{\HH B_\rho}=2^{\omega^\rho}.
\end{equation}
Let $B_0$ be a singleton; its powerspace has two points, as required.

At a successor stage, let $B_{\rho+1}$ be the bouquet of countably many
copies of $B_\rho$.  The closed union of $k$ whole components has stature
$\omega^\rho\cdot k$.  Every proper closed subset lies in such a finite
union.  It follows from the rank definition that
\[
 \st{B_{\rho+1}}
 =\sup_{k<\omega}(\omega^\rho\cdot k+1)=\omega^{\rho+1}.
\]
The powerspace of that closed $k$-component subspace embeds as a closed
subspace of $\HH B_{\rho+1}$.  By \eqref{eq:Hsum}, the external product
formula, and \eqref{eq:hom}, its stature is
\[
 \natprod_{i=1}^k2^{\omega^\rho}=2^{\omega^\rho\cdot k}.
\]
Taking the supremum over $k$ gives the lower bound
$2^{\omega^{\rho+1}}$, by continuity of exponentiation.  The upper bound
is \cref{thm:upper}, so equality holds.

At a nonzero limit $\rho$, form the bouquet of one copy of $B_\xi$ for
each $\xi<\rho$.  Every finite natural sum of the ordinals
$\omega^\xi$, $\xi<\rho$, is below $\omega^\rho$, while the individual
component statures are cofinal in $\omega^\rho$.  The rank definition
therefore gives $\st{B_\rho}=\omega^\rho$.  Each component is closed,
so the stature of its powerspace is a lower bound on
$\st{\HH B_\rho}$.  Hence
\[
 \st{\HH B_\rho}\geq\sup_{\xi<\rho}2^{\omega^\xi}
 =2^{\omega^\rho}.
\]
Again use the universal upper bound to get equality.  This proves
\eqref{eq:Btarget} at every ordinal stage.

For $\alpha>0$, expand the finite coefficients in its Cantor normal form
and write
\[
 \alpha=\omega^{\rho_1}+\cdots+\omega^{\rho_k},
 \qquad\rho_1\geq\cdots\geq\rho_k.
\]
Take $T_\alpha=B_{\rho_1}\ds\cdots\ds B_{\rho_k}$.  Its stature is the
natural sum of the summand statures, equal to the displayed Cantor
normal form.  Its Hoare powerspace is their product, whose stature is
\[
 \natprod_{i=1}^k2^{\omega^{\rho_i}}
 =2^{\natsum_{i=1}^k\omega^{\rho_i}}=2^\alpha.
\]
For $\alpha=0$ take the empty space.  Every construction preserves
$T_1$.  Recursion over countable ordinals uses only countable bouquets
and finite sums, proving the countability assertion.
\end{proof}

Combining \cref{thm:upper,prop:lowerSharp,thm:sharpT1} determines both
extreme points of $\Spec(\alpha)$ for every ordinal $\alpha$.

\section{Exactly when the older upper bound is attainable}
\label{sec:oldattain}

\begin{theorem}\label{thm:oldattain}
For an infinite ordinal $\alpha$, the following conditions are equivalent:
\begin{enumerate}[label=\textup{(\roman*)}]
\item There is a Noetherian $X$ with $\st X=\alpha$ and
      $\st{\HH X}=\omega^\alpha$.
\item $2^\alpha=\omega^\alpha$.
\item $\omega\cdot\alpha=\alpha$.
\item Every exponent in the Cantor normal form of $\alpha$ is infinite.
\item $\alpha=\omega^\omega\cdot\delta$ for some nonzero ordinal $\delta$.
\end{enumerate}
\end{theorem}
\begin{proof}
The sharp upper bound proves the equivalence of (i) and (ii), since
$2^\alpha\leq\omega^\alpha$ and a witness for $2^\alpha$ always exists.
Write $\alpha=\omega q+n$ with $n<\omega$.  Uniqueness of Cantor normal
form in
\[
 2^\alpha=\omega^q2^n
\]
shows that this equals $\omega^\alpha$ exactly when $n=0$ and $q=\alpha$.
These are precisely the conditions $\alpha=\omega\alpha$.

Write $\alpha=\omega^{\rho_1}c_1+\cdots+\omega^{\rho_k}c_k$ in Cantor
normal form.  Left multiplication by $\omega$ replaces $\rho_i$ by
$1+\rho_i$.  The normal forms agree exactly when every
$1+\rho_i=\rho_i$, that is, when every $\rho_i$ is infinite.

Finally, multiplying a Cantor normal form on the left by $\omega^\omega$
changes its exponents $\xi$ to $\omega+\xi$.  This is a strictly
increasing map whose range consists of all ordinals at least $\omega$:
for any $\rho\geq\omega$, the order type of the tail after an initial
segment of type $\omega$ supplies $\xi$ with $\rho=\omega+\xi$.
Thus precisely the normal forms in (iv) are obtained in (v).
\end{proof}

\begin{table}[htbp]
\centering
\begin{tabular}{@{}lll@{}}
\toprule
Input stature $\alpha$ & Sharp maximum $2^\alpha$ & Equals $\omega^\alpha$?\\
\midrule
$\omega$ & $\omega$ & No\\
$\omega+n$, $0<n<\omega$ & $\omega\cdot2^n$ & No\\
$\omega^2$ & $\omega^\omega$ & No\\
$\omega^3$ & $\omega^{\omega^2}$ & No\\
$\omega^\omega$ & $\omega^{\omega^\omega}$ & Yes\\
$\omega^\omega+1$ & $\omega^{\omega^\omega}\cdot2$ & No\\
$\omega^\omega+\omega$ & $\omega^{\omega^\omega+1}$ & No\\
$\omega^{\omega+1}$ & $\omega^{\omega^{\omega+1}}$ & Yes\\
$\varepsilon_0$ & $\varepsilon_0$ & Yes\\
\bottomrule
\end{tabular}
\caption{Consequences of the binary-exponent formula.  In the last row,
$\varepsilon_0$ is the least positive fixed point of $\alpha\mapsto\omega^\alpha$.}
\label{tab:ordinals}
\end{table}

This settles the upper-attainment part of the chosen question, but
extremal values alone do not classify the intervening spectrum.  The
next sections give such a classification at the first infinite ordinal
and all of its finite successors.

\section{A canonical infinite core at stature \texorpdfstring{$\omega+n$}{omega+n}}
\label{sec:core}

For a finite poset $P$, let $\JJ(P)$ be its poset of downsets, ordered
by inclusion, and let $j(P)=|\JJ(P)|$.  The empty downset is included,
so $j(\varnothing)=1$.

\begin{proposition}\label{prop:core}
Let $X$ be a $T_0$ Noetherian space with $\st X=\omega+n$, $n<\omega$.
There is a unique closed subset $F$ of stature $\omega$.  Its complement
$V$ has exactly $n$ points.  Every closed subset of $X$ is either finite
or contains $F$.  The infinite closed subsets are exactly
\begin{equation}\label{eq:infiniteclosed}
 F\cup D,\qquad D\in\JJ(Q),
\end{equation}
where $Q$ is the specialization poset on $V$.
\end{proposition}
\begin{proof}
For $n=0$ take $F=X$.  For $n>0$ use \cref{cor:tail}.  Every proper
closed subset of $F$ has stature less than $\omega$, and hence is finite
by \cref{lem:finite}.  Also $F$ is infinite.

If a closed $C\subseteq X$ does not contain $F$, then $C\cap F$ is a
proper closed subset of $F$, hence finite.  Its remaining points lie
in the finite set $V$, so $C$ is finite.  Conversely, a closed set
containing $F$ is infinite.  The interval $[F,X]$ is isomorphic to the
closed-set lattice of $V$.  Since $V$ is finite and $T_0$, its closed
sets are exactly the downsets of $Q$, giving \eqref{eq:infiniteclosed}.

Any other closed subset of stature $\omega$ must be infinite and hence
contain $F$.  If it strictly contained $F$, its rank would exceed
$r_X(F)=\omega$.  This proves uniqueness.
\end{proof}

Define the distinguished downset
\begin{equation}\label{eq:U}
 U=\{v\in V:\cl_X\{v\}\text{ is finite}\}.
\end{equation}
It is a downset of $Q$, because $v'\leq v$ implies
$\cl_X\{v'\}\subseteq\cl_X\{v\}$.  If $v\notin U$, its closure is
infinite and therefore contains all of $F$.  Thus $U$ distinguishes the
finite-closure points from those whose closure absorbs the whole core.

\subsection{The generic-point issue}

A point $g\in F$ can have closure all of $F$, even though every proper
closed subset of $F$ is finite.  Such a point would make $F$ itself
finitarily generated, changing the apparent list of points of $\Hfin X$.
It must not be silently assumed absent.

\begin{lemma}[Generic-point normalization]\label{lem:generic}
There is at most one point $g\in F$ with $\cl_X\{g\}=F$.  If it exists,
put $X'=X\setminus\{g\}$.  The trace map
$\Cl(X)\to\Cl(X')$ is an order isomorphism.  Removing $g$ preserves
$\st X$, $\st{\HH X}$, $Q$, and $U$.  After this removal, every point of
the remaining core has finite closure.
\end{lemma}
\begin{proof}
Two points with closure $F$ would be indistinguishable, contrary to
$T_0$.  Every finite closed set avoids $g$, since any closed set
containing it contains the infinite set $F$.  Each infinite closed set
has the unique form $F\cup D$, and its trace is
$(F\setminus\{g\})\cup D$.

Taking these traces preserves and reflects inclusion.  For two finite
closed sets nothing changes.  For two infinite ones inclusion is
controlled by their $D$ parts.  A trace of an infinite set cannot be
contained in a finite set, since $F\setminus\{g\}$ is still infinite.
For a finite set contained in the trace of an infinite one, the same
inclusion held before removal.  The trace map is onto the closed sets
of the subspace by definition, hence is an order isomorphism.

Stature is therefore preserved, and the powerspaces are homeomorphic
by \cref{prop:closedH}.  The specialization order restricted to $V$ is
unchanged.  An infinite closure loses at most $g$ and stays infinite;
a finite closure never contained $g$.  Thus $U$ is unchanged.  Every
other point of $F$ already had a proper, therefore finite, closure.
\end{proof}

For the rest of the spectrum argument, we perform this normalization
and again call the resulting space $X$ and its core $F$.  This is a
closed-lattice simplification, not a claim that the original space
lacked a generic point.

\section{A finite-cell rank lemma}
\label{sec:cells}

The remaining transfinite computation can be packaged in an elementary
lemma.  It explains why a finite downset count becomes the coefficient
of $\omega$.

\begin{lemma}[Finite-cell rank]\label{lem:cells}
Suppose a $T_0$ space $Y$ is partitioned into infinite sets $A_p$ indexed
by a finite poset $P$, and a finite exceptional set $E$.  Assume:
\begin{enumerate}[label=\textup{(\alph*)}]
\item Every closed $C\subseteq Y$ consists of all the cells indexed by
      some downset $I(C)\subseteq P$, together with finitely many other
      points.
\item For every downset $I\subseteq P$, the pure union
      $A_I=\bigcup_{p\in I}A_p$ is closed.
\item Every point belonging to a cell $A_p$ has finite closure in $Y$.
\end{enumerate}
Then $Y$ is Noetherian.  For every closed $C$,
\begin{equation}\label{eq:cellrank}
 r_Y(C)=\omega\cdot|I(C)|+|C\setminus A_{I(C)}|.
\end{equation}
In particular,
\begin{equation}\label{eq:celltop}
 \st Y=\omega\cdot|P|+|E|.
\end{equation}
\end{lemma}
\begin{proof}
The set $I(C)$ is uniquely determined: a cell is infinite, so a finite
remainder cannot contain an additional entire cell.  Define $w(C)$ to
be the right side of \eqref{eq:cellrank}.

If $C\subsetneq C'$, then $I(C)\subseteq I(C')$.  If the index sets
are different, the coefficient of $\omega$ strictly increases.  If
they agree, the finite remainder strictly increases in cardinality.
Thus $w(C)<w(C')$ in either case.  No infinite descending chain of
closed sets is possible, and induction gives $r_Y(C)\leq w(C)$.

For the reverse inequality, induct on $k=|I(C)|$.  When $k=0$, $C$ is
a finite $T_0$ space, so its stature is its cardinality.
Let $k>0$ and first consider the pure union $A_I$, $|I|=k$.
Choose a maximal $p\in I$ and put $I'=I\setminus\{p\}$, again a downset.
For any positive integer $N$, choose $N$ distinct points in the infinite
cell $A_p$.  The union $K_N$ of their closures is finite by (c).
It lies inside $A_I$, since $A_I$ is closed.  The closed set
\[
 B_N=A_{I'}\cup K_N
\]
is a proper subset of $A_I$, has exactly $k-1$ full cells, and has at
least $N$ points outside those cells.  The induction hypothesis gives
\[
 r_Y(B_N)\geq\omega\cdot(k-1)+N.
\]
Taking the supremum over $N$ in the rank recursion for $A_I$ yields
$r_Y(A_I)\geq\omega\cdot k$.

Now return to an arbitrary closed $C$ with $I(C)=I$.  The complement
$C\setminus A_I$ has some finite cardinality $m$.  Apply the lower
closed--open inequality inside $C$, with closed part $A_I$.  Its finite
$T_0$ complement has stature $m$, so
\[
 r_Y(C)=\st C\geq\st{A_I}+m\geq\omega\cdot k+m=w(C).
\]
This matches the upper bound.  Taking $C=Y$ proves \eqref{eq:celltop}.
\end{proof}

It is essential that cell points have finite closure and that every
pure downset union is closed.  Merely partitioning a space into finitely
many infinite pieces would not determine its stature.

\section{The complete spectrum at \texorpdfstring{$\omega+n$}{omega+n}}
\label{sec:spectrum}

\subsection{The formula for an arbitrary space}

\begin{theorem}\label{thm:formula}
Let $X$ be a Noetherian space of stature $\omega+n$, $n<\omega$.
Pass to $T_0$, let $F$ be the core of \cref{prop:core}, let $Q$ be the
$n$-point complement poset, and let $U$ be \eqref{eq:U}.  Then
\begin{equation}\label{eq:spectrumformula}
 \st{\HH X}=\omega\cdot j(U)+(j(Q)-j(U)).
\end{equation}
\end{theorem}
\begin{proof}
Use \cref{lem:generic} to remove a generic core point if necessary.
Every point of $F$, and every point of $U$, now has finite closure.
A finite union of their closures is finite.  A finite generating set
has infinite closure precisely when it contains some point of $V\setminus U$.

It follows that the points of $\Hfin X$ are all finite closed sets,
together with precisely the following infinite closed sets:
\begin{equation}\label{eq:exceptional}
 E=\{F\cup D:D\in\JJ(Q),\ D\nsubseteq U\}.
\end{equation}
For the converse implicit in this assertion, if $D$ meets
$V\setminus U$, then the closure of the finite set $D$ is exactly
$F\cup D$.  Indeed, its trace on $V$ is the downset $D$, and one of its
point closures contains $F$.  Every finite closed set is trivially
finitarily generated by all of its own points.

For $J\in\JJ(U)$, define a cell
\begin{equation}\label{eq:AJ}
 A_J=\{C\in\Cl(X): C\text{ is finite and }C\cap V=J\}.
\end{equation}
A downset of $U$ is also a downset of $Q$, because $U$ itself is a
downset.  Every finite closed set belongs to exactly one such cell.

Each $A_J$ is infinite.  The finite closed set
$B_J=\bigcup_{v\in J}\cl_X\{v\}$ has trace $J$ on $V$.  There are
infinitely many different finite closed sets contained in $F$:
otherwise their finite union would contain every point of the infinite
core, because each such point has finite closure.  More directly,
one can repeatedly choose a core point outside the current finite
closed union and adjoin its closure.  Union with $B_J$ therefore
produces infinitely many distinct members of $A_J$.

The closure in $\Hfin X$ of a point $C\in A_J$ consists of finitarily
closed subsets contained in the finite set $C$.  It is finite.
Thus condition (c) of \cref{lem:cells} holds.

We next verify its closed-set conditions, with index poset $P=\JJ(U)$.
The trace of $\boxF C$ is finite when $C$ is finite.  When $C=F\cup D$
is infinite, its trace contains exactly the entire cells $A_J$ with
$J\subseteq D\cap U$, and some of the finitely many exceptional points
in \eqref{eq:exceptional}.  By \cref{prop:closedH}, every closed subset
of $\Hfin X$ is a finite union of such traces.  Its full-cell indices
are therefore a downset of $\JJ(U)$; all other points form a finite
remainder.  This is condition (a).

Conversely, if $I$ is a downset of $\JJ(U)$, take the union of the traces
of $\boxF{F\cup J}$ over the maximal $J\in I$.  Each $J$ is contained
in $U$, so this trace has no exceptional point and consists exactly
of the cells with index contained in $J$.  The union is precisely
$\bigcup_{J\in I}A_J$.  It is closed, proving condition (b).  For
$I=\varnothing$ use the empty union.

The finite-cell lemma now gives
\[
 \st{\Hfin X}=\omega\cdot|\JJ(U)|+|E|.
\]
The downsets of $Q$ contained in $U$ are exactly the downsets of $U$,
so $|E|=j(Q)-j(U)$.  Finally use \cref{prop:finitary}.
\end{proof}

\subsection{Every finite pair has a countable realization}

\begin{theorem}\label{thm:realization}
For every finite poset $Q$ of size $n$ and every downset $U\subseteq Q$,
there is a countable $T_0$ Noetherian space $X(Q,U)$ with
\[
 \st{X(Q,U)}=\omega+n,
 \qquad
 \st{\HH X(Q,U)}=\omega\cdot j(U)+(j(Q)-j(U)).
\]
\end{theorem}
\begin{proof}
Let $C$ be a countably infinite set disjoint from $Q$.  On the underlying
set $C\ds Q$, declare the closed subsets to be exactly the following:
\begin{align}
 & A\ds D, && A\subseteq C\text{ finite},\quad D\in\JJ(U);
       \label{eq:modelFinite}\\
 & C\ds D, && D\in\JJ(Q). \label{eq:modelInfinite}
\end{align}
The empty set and the whole space occur.  Finite unions preserve these
forms, since unions of downsets are downsets.  For an arbitrary
intersection, if every member has form \eqref{eq:modelInfinite}, so
does the intersection.  Otherwise fix a finite-type member; the
intersection has finite $C$ part and a downset contained in $U$, hence
has form \eqref{eq:modelFinite}.  This is a topology.

A descending chain of infinite-type closed sets can change only its
finite $Q$ part, hence can have only finitely many strict steps before
it either stops or enters a finite-type closed set.  Thereafter it lies
inside a finite set.  Thus the topology is Noetherian.

The subspace $C$ is closed and has the cofinite topology, of stature
$\omega$: its proper closed subsets are all finite sets.  The complement
is exactly the finite $T_0$ space $Q$, of stature $n$.  Both sides of
the closed--open inequality equal $\omega+n$, so
$\st{X(Q,U)}=\omega+n$.

For $v\in U$, its closure is $\mathord\downarrow_Q v$, a finite set.
For $v\notin U$, its closure is
$C\ds\mathord\downarrow_Q v$.  Points of $C$ have singleton closures.
These closures are pairwise distinct, so the space is $T_0$.
They also show that the distinguished finite-closure downset is exactly
the specified $U$.  Apply \cref{thm:formula}.
\end{proof}

\begin{corollary}[Exact finite-tail spectrum]\label{cor:spectrum}
For every $n<\omega$,
\[
 \Spec(\omega+n)=
 \{\omega\cdot j(U)+(j(Q)-j(U)):
                    |Q|=n,\ U\in\JJ(Q)\}.
\]
Every value has a countable witness.
\end{corollary}

\begin{corollary}[Rigidity in the $T_1$ case]\label{cor:T1tail}
If $X$ is $T_1$ and $\st X=\omega+n$, then
\[
 \st{\HH X}=\omega\cdot2^n.
\]
In fact $X$ is the disjoint sum of an infinite cofinite space and an
$n$-point discrete space.
\end{corollary}
\begin{proof}
The finite complement of the core is closed because $X$ is $T_1$.
Thus the core and its complement are both clopen.  The core has all
finite subsets closed and no other proper closed subsets, hence is
cofinite.  Its finite complement is discrete.  Equivalently, in the
spectrum formula $Q$ is an antichain and $U=Q$, giving $j(U)=2^n$ and
zero finite remainder.
\end{proof}

The source space can have a complicated family of finite point closures
inside its core.  At stature $\omega+n$, none of that complexity changes
the answer: after the justified normalization, only the finite pair
$(Q,U)$ survives in the formula.

\section{Examples, gaps, and exhaustive finite data}
\label{sec:examples}

\subsection{The first spectra}

For finite input $n$, a $T_0$ space is an $n$-point poset, and its
powerspace has one point per downset.  Therefore the finite-input
spectrum is simply
\[
 \Spec(n)=\{j(Q):|Q|=n\}.
\]
For infinite input with finite tail, the first three spectra are
\begin{align*}
 \Spec(\omega)&=\{\omega\},\\
 \Spec(\omega+1)&=\{\omega+1,\ \omega\cdot2\},\\
 \Spec(\omega+2)&=\{\omega+2,\ \omega+3,\ \omega\cdot2+1,\
                   \omega\cdot2+2,\ \omega\cdot3,\ \omega\cdot4\}.
\end{align*}

For $n=1$, $Q$ has two downsets.  Choosing $U=\varnothing$ gives
$\omega+1$; choosing $U=Q$ gives $\omega\cdot2$.  Thus, for example,
$\omega+2$ lies between the extreme values but is not itself attainable
at input $\omega+1$.

For $n=2$, there are two poset types.  A chain has $j(Q)=3$ and possible
$j(U)=1,2,3$, giving $\omega+2$, $\omega\cdot2+1$, and $\omega\cdot3$.
An antichain has $j(Q)=4$ and possible $j(U)=1,2,4$, giving
$\omega+3$, $\omega\cdot2+2$, and $\omega\cdot4$.

For $n=3$, the eighteen values can be presented without writing a long
ordinal list: for each coefficient $a$ in \cref{tab:n3}, take
$\omega\cdot a+b$ for each finite tail $b$ in the second column.

\begin{table}[htbp]
\centering
\begin{tabular}{@{}cl@{}}
\toprule
Coefficient $a$ & Permitted tails $b$ at input $\omega+3$\\
\midrule
$1$ & $3,4,5,7$\\
$2$ & $2,3,4,6$\\
$3$ & $1,2,3$\\
$4$ & $0,1,2,4$\\
$5$ & $0$\\
$6$ & $0$\\
$8$ & $0$\\
\bottomrule
\end{tabular}
\caption{The complete spectrum at $\omega+3$, encoded as pairs $(a,b)$.}
\label{tab:n3}
\end{table}

Both endpoints of every finite-tail spectrum have elementary witnesses.
Take $Q$ to be a chain and $U=\varnothing$ to get
$\omega+(j(Q)-1)=\omega+n$.  Take $Q$ to be an antichain and $U=Q$ to get
$\omega\cdot2^n$.  The intervening gaps are restrictions on finite
poset downset counts, not failures of ordinal continuity.

\subsection{Enumeration and certificates}

The accompanying \texttt{verify\_spectra.py} uses only the Python standard
library.  For each $n$, it enumerates all subsets of the pairs $(i,j)$
with $0\leq i<j<n$, retaining exactly the transitive relations.
Every finite poset admits a linear extension, so relabeling along one
puts it into this collection.  This covers every isomorphism type,
usually more than once.  The counts are \emph{not} counts of unlabeled
posets and \emph{not} counts of all labeled posets.

For each retained relation it enumerates all downsets $U$, counts
$j(Q)$ and $j(U)$, and stores the pair
$(j(U),j(Q)-j(U))$.  A concrete relation and a distinguished downset
are retained as a witness for each distinct pair.

\begin{table}[htbp]
\centering
\begin{tabular}{@{}rrr@{}}
\toprule
$n$ & Naturally labeled transitive relations & $|\Spec(\omega+n)|$\\
\midrule
$0$ & $1$ & $1$\\
$1$ & $1$ & $2$\\
$2$ & $2$ & $6$\\
$3$ & $7$ & $18$\\
$4$ & $40$ & $53$\\
$5$ & $357$ & $154$\\
$6$ & $4824$ & $431$\\
\bottomrule
\end{tabular}
\caption{Exhaustive finite enumeration performed for this manuscript.}
\label{tab:counts}
\end{table}

Across these runs the program examines $101{,}659$ pairs consisting of
a retained poset and one of its downsets.  The archive contains the
full spectra in \texttt{spectra.csv}, one witness per distinct value in
\texttt{witnesses.json}, and the run summary in
\texttt{verification\_results.json}.  The spectrum-size sequence
$1,2,6,18,53,154,431$ is reported as computed data, not asserted to be
new to integer-sequence databases.

\subsection{An independent finite-model identity}

There is a useful check of the realization formula that does not use
ordinal ranks.  Replace the infinite cofinite core by an $N$-point
antichain $C_N$.  Keep the order on $Q$, put each point of $C_N$ below
each point of $Q\setminus U$, and add no further cross-comparabilities.
Because $U$ is a downset, $Q\setminus U$ is an upset and this is a
partial order.  Call it $P_{N,Q,U}$.

\begin{proposition}\label{prop:finitecheck}
For $N\geq1$,
\[
 j(P_{N,Q,U})=2^N j(U)+(j(Q)-j(U)).
\]
\end{proposition}
\begin{proof}
A downset whose $Q$ part is contained in $U$ has any of the $2^N$
possible core parts, and its $Q$ part has $j(U)$ choices.  A downset
whose $Q$ part is not contained in $U$ must contain all of $C_N$,
and it has $j(Q)-j(U)$ choices for its $Q$ part.  These cases are
disjoint and exhaustive.
\end{proof}

For every retained $Q$ with at most four points, every downset $U$, and
each $N=1,2,3$, the program builds this finite poset and independently
enumerates its downsets.  All $1{,}221$ checks matched the formula.
This is an arithmetic and construction check, \emph{not} a computer
proof of \cref{thm:formula}.  In particular, taking an ordinary supremum
of the finite counts would lose the ordinal coefficient information;
the finite-cell lemma is what justifies the transfinite rank.

\section{Dependency audit and the boundary of the result}
\label{sec:audit}

The argument has two largely independent parts.  The universal bound
uses finite unions, the closed--open rank inequality, and the fact that
$2^\lambda$ is naturally additively indecomposable at a nonzero limit
$\lambda$.  The finite-tail spectrum uses the unique core, the finitary
closed-lattice isomorphism, and the finite-cell rank computation.
Their overlap is mostly notation and elementary rank monotonicity.

Several potential errors deserve explicit attention.

\paragraph{Ordinal, not cardinal, arithmetic.}
The expression $2^\omega$ here is $\omega$.  Also, the lower bound is
$1+\alpha$, not $\alpha+1$.  At $\alpha=\omega$ the former is $\omega$,
and replacing it by the latter would contradict the exact spectrum.
Natural sums occur in finite disjoint sums and natural products in the
imported product theorem; ordinary sums occur in a closed--open lower
bound and in the final normal form $\omega a+b$.

\paragraph{The empty powerspace point.}
Even when $X$ is empty, $\HH X$ has the one point $\varnothing$.
For a finite $n$-point $T_0$ space, the stature of $\HH X$ is its number
of closed subsets, not that number minus one.  These conventions fix
the base cases and the finite remainders in the spectrum formula.

\paragraph{No hidden assumption that all point closures are finite.}
At stature $\omega+n$, a core generic point may exist.  It is removed
only after proving an isomorphism of closed lattices.  A point outside
$U$ is never removed: its infinite closure is exactly what produces
one of the exceptional points in the finitary powerspace.

\paragraph{The product dependency is localized.}
Equation \eqref{eq:productexternal} is a published theorem, not reproved
in this manuscript.  It is used to certify the lower bounds for the
$T_1$ bouquet witnesses.  The universal upper proof, the old-bound
nonattainment implications, and the complete $\omega+n$ classification
can be read without it.  Existing wqo sharpness witnesses give a
separate route to unrestricted upper-bound attainment.

\paragraph{What has and has not been solved.}
The manuscript supplies a proof of a sharper general bound, witnesses
for both extreme points at every ordinal, an exact answer to the
upper-attainment subquestion, and the full spectrum for $\omega+n$.
It does not classify $\Spec(\omega\cdot2)$, $\Spec(\omega^2)$, or all
higher spectra.  There is no claim that a spectrum is an interval;
the first nontrivial finite tails already show otherwise.

\paragraph{Verification status.}
The displayed transfinite arguments are mathematical proofs proposed
here and have been checked against the stated definitions during
preparation.  They have not received independent human peer review
and have not been formalized in Lean or another proof assistant.  The
included executable checks concern only finite posets and the finite
realization identity.  The novelty search located the decisive earlier
wqo theorem and the original Noetherian question, but does not establish
priority for the general extension or the spectrum theorem.

\subsection{Specific next mathematical questions}

The next natural target is stature $\omega\cdot2$.  The simple core
argument no longer makes all proper closed subsets finite: some can
already have stature $\omega$, and their powerspace pieces have
transfinite internal structure.  A promising replacement for the
finite-cell lemma would attach ordinal weights, rather than finite
cardinalities, to a finite poset of closed strata.  One would need
precise hypotheses ensuring that finite unions of fragments cannot
silently complete a stratum.

A second question is purely finite: describe the set of pairs
$(j(U),j(Q)-j(U))$ for $n$-point $Q$ without enumerating all posets.
The spectrum theorem reduces the first infinite layer completely to
that finite combinatorial problem.  Bounds on possible downset counts,
structural recurrences for ordinal sums and disjoint sums of finite
posets, and efficient isomorph rejection could all improve the
computation without touching the transfinite proof.

A third question concerns the $T_1$ subclass at higher statures.  At
$\omega+n$ it is rigid and always attains the maximum.  The bouquet
construction provides maximal $T_1$ examples everywhere, but it does
not prove that every $T_1$ space at every ordinal stature is maximal.
No such extrapolation is used in this article.

\appendix
\section{Reproducing the artifacts}

The archive contains this article as \texttt{.tex} and \texttt{.pdf}, a
reusable Bib\TeX{} file, the Python program, computed spectra, explicit
finite witnesses, and proof-audit and literature notes.  Third-party
source PDFs and font files are not included.

With Python~3.9 or later, run
\begin{quote}\ttfamily
python verify\_spectra.py
\end{quote}
The defaults cover $0\leq n\leq6$ and finite-model checks through $n=4$.
Options are \texttt{--max-n}, \texttt{--finite-check-n}, and
\texttt{--output}; the input guard allows at most $n=7$ because the
enumeration is exponential.  Details and file descriptions are in
\texttt{README.md}.

With a normal TeX Live installation, compile using
\begin{quote}\ttfamily
latexmk -pdf -interaction=nonstopmode -halt-on-error article.tex
\end{quote}
The bibliography is embedded in the TeX source, so a separate Bib\TeX{}
run is unnecessary.

\begin{thebibliography}{9}
\bibitem{GL23}
Jean Goubault-Larrecq and Bastien Laboureix.
\newblock \emph{Statures and sobrification ranks of Noetherian spaces}.
\newblock Houston Journal of Mathematics \textbf{49}(1) (2023), 1--76.
\newblock Public version consulted: arXiv:2112.06828v5, submitted
24 February 2023; 68-page preprint, with Question~1 on printed page~63.
\newblock \href{https://arxiv.org/abs/2112.06828v5}{arXiv:2112.06828v5}.

\bibitem{AHLLSSV24}
Sergio Abriola, Simon Halfon, Aliaume Lopez, Sylvain Schmitz,
Philippe Schnoebelen, and Isa Vialard.
\newblock \emph{Measuring well-quasi-ordered finitary powersets}.
\newblock arXiv:2312.14587v2, 17 July 2024.
\newblock Version consulted for Theorems~3.3 and~4.3.
\newblock \href{https://arxiv.org/abs/2312.14587v2}{arXiv:2312.14587v2}.
\end{thebibliography}

\end{document}
